\section*{Research Foundations and Recommended Reading}

The chapter is an original synthesis of microeconomics, operations, service management, learning and organization studies, and contemporary AI-productivity research. Its mathematical models describe \textbf{specific assumptions}, and its Bayes Institute models are illustrative. These sources should be used to sharpen a founder's reasoning, not to assign universal numerical performance improvements to an untested company.

\paragraph*{E1: William J. Baumol (1967), ``Macroeconomics of Unbalanced Growth: The Anatomy of Urban Crisis,'' \emph{American Economic Review}, 57(3), 415 to 426.} A foundation for understanding why labor-intensive services such as education may experience relative cost pressure when productivity grows faster elsewhere. It does not claim educational quality cannot improve or that productivity in all educational tasks is static.

\paragraph*{E2: David J. Deming, Claudia Goldin, Lawrence F. Katz, and Noam Yuchtman (2015), ``Can Online Learning Bend the Higher Education Cost Curve?'' \emph{American Economic Review}, 105(5), 496 to 501.} Finds evidence consistent with lower online-education prices in a studied historical period but notes uncertainty about quality. It is not evidence that every online program is cheaper per verified learning gain.

\paragraph*{E3: John D. C. Little (1961), ``A Proof for the Queuing Formula: $L = \lambda W$,'' \emph{Operations Research}, 9(3), 383 to 387.} The fundamental average throughput, work-in-process, and time relationship under the relevant stability assumptions.

\paragraph*{E4: T. P. Wright (1936), ``Factors Affecting the Cost of Airplanes,'' \emph{Journal of the Aeronautical Sciences}, 3(4), 122 to 128.} An early account of production learning curves, useful as a conceptual antecedent but not a numerical prediction for faculty work.

\paragraph*{E5: W. Edwards Deming (1986), \emph{Out of the Crisis}.} Explains how quality, variation, management systems, and process improvement interact with productivity and service delivery.

\paragraph*{E6: James L. Heskett, Thomas O. Jones, Gary W. Loveman, W. Earl Sasser Jr., and Leonard A. Schlesinger (1994), ``Putting the Service-Profit Chain to Work,'' \emph{Harvard Business Review}, March to April.} A management framework linking internal service, employees, customer value, and performance; its causal links must be tested rather than assumed in any particular firm.

\paragraph*{E7: G. Lynn Shostack (1984), ``Designing Services That Deliver,'' \emph{Harvard Business Review}, January.} A foundational approach to mapping visible service and backstage processes.

\paragraph*{E8: Erik Brynjolfsson, Danielle Li, and Lindsey R. Raymond (2025), ``Generative AI at Work,'' \emph{Quarterly Journal of Economics}, 140(2), 889 to 942.} Field evidence of heterogeneous productivity effects from AI assistance in a customer-support setting. The average effect is not a general educational labor-saving coefficient.

\paragraph*{E9: Shakked Noy and Whitney Zhang (2023), ``Experimental Evidence on the Productivity Effects of Generative Artificial Intelligence,'' \emph{Science}, 381(6654), 187 to 192.} Randomized evidence for selected professional writing tasks, with substantial observed effects on task completion and output quality in that setting.

\paragraph*{E10: Richard E. Mayer and Logan Fiorella, eds. (2021), \emph{The Cambridge Handbook of Multimedia Learning}, 3rd ed.} A research-oriented reference on the design of educational media and management of cognitive processing. Use alongside Chapters 11 to 13 to test whether lower-cost explanations remain pedagogically sound.

\paragraph*{E11: Susan Athey and Scott Stern, ``An Empirical Framework for Testing Theories About Complementarity in Organizational Design.''} An econometric examination of the difficulty of identifying complementary organizational practices and related selection bias; useful for resisting unsupported synergy claims.

\paragraph*{E12: Amy C. Edmondson (1999), ``Psychological Safety and Learning Behavior in Work Teams,'' \emph{Administrative Science Quarterly}, 44(2), 350 to 383.} Research connecting team psychological safety to learning behavior in a studied organizational setting. Relevant to building a teaching and operations culture where errors are surfaced and repaired, not concealed to protect dashboards.

\paragraph*{E13: Michael E. Porter (1996), ``What Is Strategy?'' \emph{Harvard Business Review}.} Activity-system fit and trade-offs, useful for understanding why a differentiated service must be supported by a coherent cost and operating design.

\paragraph*{E14: Robert J. Mislevy, Linda S. Steinberg, and Russell G. Almond (2003), ``On the Structure of Educational Assessments.''} A foundation for distinguishing inexpensive scoring from valid evidence of capability. See Chapter 13 for detailed treatment.

\paragraph*{E15: National Academies of Sciences, Engineering, and Medicine (2018), \emph{How People Learn II: Learners, Contexts, and Cultures}.} A synthesis reminding institutions that capability depends on prior knowledge, learning conditions, motivation, and context, not only resource efficiency.

\subsection*{How to Read the Apparent Tensions}

Baumol illuminates persistent labor intensity; Wright illuminates productive learning through repetition; Little illuminates flow and delay; Deming illuminates systemic quality costs; service-blueprinting and service-profit-chain theories illuminate operations and relationships; contemporary AI studies show that some cognitive tasks can become substantially more productive. These are \textbf{different mechanisms}, not mutually exclusive verdicts about education.

An institute can experience rising costs for senior expert attention while falling costs for content distribution. It can become more expensive per learner and cheaper per verified outcome. It can automate feedback drafting and discover that the hard work of assessing independent judgment remains. It can earn more revenue by increasing cohort size while simultaneously reducing educational quality. The founder must examine which of these mechanisms is operating in their own system and at what margin.

The 1 metric to distrust most is the metric whose improvement would be celebrated \textbf{even if the customer had become worse off}.
